\chapter{Measuring Algorithms}\label{ch:measuring}

\begin{center}\itshape
Beware of bugs in the above code; I have only proved it correct, not tried it.\\
\upshape(Donald Knuth, end of a memo on van Emde Boas's priority queues, 1977)
\end{center}

In 1844 Gabriel Lam\'e had just been elected to the Paris Academy of Sciences. As a young engineer he had spent twelve years in
St Petersburg, teaching at the Institute of Ways of Communication and helping to plan the roads and bridges of the city, and he would
be remembered for the equations of elasticity. That year he sent the Academy a three-page note on a question nobody had thought worth asking for two thousand years: how many
divisions does Euclid's algorithm need? His answer was that it never needs more than five times the number of decimal digits of the
smaller number, and that the worst inputs are consecutive Fibonacci numbers (Tome~1, Theorem~1.8). Others had found weaker
bounds before him, but Lam\'e's note is often called the first paper in the analysis of algorithms: the first time someone counted
the steps of a procedure as a function of the size of its input, and proved the count sharp.

The question has become an industry. A program that is correct but needs a billion years is of no use, and the difference between a
good algorithm and a bad one is rarely a factor of two: it is the difference between $n^2$ and $n\log n$, between $2^n$ and $n^3$. On a
list of a million names, $n^2$ is $10^{12}$ steps, about a quarter of an hour at a billion steps per second, while $n\log_2 n$ is twenty
million, a fiftieth of a second. Faster computers do not close that gap; better mathematics does.

This chapter sets up the tools for the rest of the book. We fix what a ``step'' is, and the language of $O$, $\Omega$ and $\Theta$ for
comparing growth. We solve the recurrences that come from cutting a problem into pieces, once and for all, in the \emph{master theorem}.
We learn to charge an expensive operation to the cheap ones that came before it (\emph{amortised analysis}), and we use it to make
the Horn-clause solver of Tome~C1 run in linear time. And then we turn the question around: not how fast an algorithm \emph{is}, but how
fast \emph{any} algorithm \emph{can} be. Sorting $n$ items needs about $n\log_2 n$ comparisons; finding the largest and the smallest
needs exactly $\lceil 3n/2\rceil-2$; finding the best and the second best player of a tennis tournament needs
$n+\lceil\log_2 n\rceil-2$ matches, a fact first worried about by Lewis Carroll.

\begin{think}
You have $8$ coins, all of the same weight except one, which is heavier, and a two-pan balance. How many weighings do you need, in the
worst case, to find the heavy coin? Can you prove that nobody can do better? What if the odd coin may be heavier \emph{or} lighter, and
you must also say which? A weighing has three outcomes, and that is the whole idea of Section~\ref{sec:lower}.
\end{think}

\section{What a step is}\label{sec:models}

To count steps we need a machine. Real computers differ in a thousand details, but they share a skeleton, which we take as our model.

\begin{definition}\label{def:wordram}
The \emph{word RAM} has an unbounded array of memory cells, each holding a \emph{word} of $w$ bits, that is, an integer in
$\{0,1,\dots,2^w-1\}$. In one \emph{step} it can read or write a cell whose address is held in a word, and perform one of the
operations $+$, $-$, $\times$, integer division, comparison, and the bitwise operations on words. On an input of size $n$
(that is, $n$ words), the word size satisfies $w\ge\log_2 n$, so that an address into the input fits in a word.
\end{definition}

The condition $w\ge\log_2 n$ is natural, since the machine must be able to point at its own input. The other half of the bargain is
that $w$ is not \emph{too} large: in this book $w=O(\log n)$ unless we say otherwise, and an algorithm that handles bigger numbers must
pay for them, a word at a time.

\begin{think}
Why must $w$ be bounded? Starting from $x=2$ and squaring $n$ times gives $2^{2^n}$, a number with $2^n+1$ bits, in $n$ steps. If
a machine could multiply numbers of any size in one step, what else could it compute ``fast''? (Hartmanis and Simon showed in 1974
that such a machine can solve in polynomial time every problem that a normal machine solves in polynomial \emph{space}, which is
believed to be much more; Tome~C3 explains the difference.)
\end{think}

When we count steps we are rarely interested in exact counts, which depend on the machine; we want the \emph{growth}.

\begin{definition}\label{def:bigo}
Let $f,g$ be functions from the positive integers (or reals) to the reals, with $g(n)>0$ for large $n$. We write
\begin{itemize}
\item $f=O(g)$ if there are $C>0$ and $n_0$ with $|f(n)|\le Cg(n)$ for all $n\ge n_0$;
\item $f=\Omega(g)$ if there are $c>0$ and $n_0$ with $f(n)\ge cg(n)$ for all $n\ge n_0$;
\item $f=\Theta(g)$ if both hold, and $f=o(g)$ if $f(n)/g(n)\to0$.
\end{itemize}
\end{definition}

The equals sign is a convention, not an equality: $n=O(n^2)$ but not $n^2=O(n)$. The rules of computation all follow from the
definitions in one line: $O(f)+O(g)=O(\max(f,g))$, $O(f)\,O(g)=O(fg)$, and $f=O(g)$, $g=O(h)$ give $f=O(h)$ (multiply the constants
and take the larger $n_0$). Logarithms to different bases differ by a constant factor, $\log_a n=\log_b n/\log_b a$, so inside an $O$ we
do not write the base. The scale to remember is
\[
1\ll\log n\ll\sqrt n\ll n\ll n\log n\ll n^2\ll n^3\ll 2^n\ll n!\ll n^n,
\]
where $f\ll g$ means $f=o(g)$. For instance $\log n=o(n^\varepsilon)$ for every $\varepsilon>0$: with $m=\varepsilon\ln n$, the ratio $\ln n/n^\varepsilon$
is $m/(\varepsilon e^m)\le m/(\varepsilon(1+m+m^2/2))\to0$.

For an algorithm we look at the \emph{worst case}: $T(n)$ is the largest number of steps over inputs of size $n$. A bound $T(n)=O(g(n))$
is then a guarantee on every input. The \emph{average case}, over a probability distribution on inputs, is the subject of Chapter~8,
where the randomness is put inside the algorithm instead.

\subsection*{Euclid, counted in bits}

Lam\'e counted divisions. On numbers too big for a word, a division is not one step, and we should count bit operations. The schoolbook
method divides an $m$-bit number by an $\ell$-bit number, $m\ge\ell$, by $m-\ell+1$ subtractions of $(\ell+1)$-bit numbers, one per bit of the
quotient: about $C\ell(m-\ell+1)$ bit operations for a constant $C$. Write $\ell(x)=\lfloor\log_2x\rfloor+1$ for the number of bits of $x\ge1$.

\begin{proposition}\label{prop:euclidbits}
With schoolbook division, Euclid's algorithm on $a>b\ge1$ uses $O(\ell(a)\,\ell(b))$ bit operations.
\end{proposition}
\begin{proof}
Write $r_0=a$, $r_1=b$ and $r_{i-1}=q_ir_i+r_{i+1}$ with $0\le r_{i+1}<r_i$, for $i=1,\dots,k$, until $r_{k+1}=0$. Since $r_{i-1}<(q_i+1)r_i$,
\[
\ell(r_{i-1})\le\ell(q_i+1)+\ell(r_i)\le\ell(q_i)+1+\ell(r_i),
\]
so the $i$-th division costs at most $C\ell(r_i)(\ell(q_i)+2)\le C\ell(b)(\log_2q_i+3)$. Next, $r_{i-1}\ge q_ir_i$, so
$q_1q_2\cdots q_k\le\frac{r_0}{r_1}\frac{r_1}{r_2}\cdots\frac{r_{k-1}}{r_k}=\frac{a}{r_k}\le a$. Summing,
\[
\text{cost}\le C\ell(b)\Bigl(\log_2\prod_{i=1}^kq_i+3k\Bigr)\le C\ell(b)\bigl(\log_2a+3k\bigr).
\]
By Lam\'e's theorem, $b\ge F_{k+1}\ge\varphi^{k-1}$ with $\varphi=(1+\sqrt5)/2$ (Tome~1, Theorem~1.8 and Exercise~1.6), so
$k\le1+\log_\varphi b=O(\ell(b))=O(\ell(a))$, and the cost is $O(\ell(a)\ell(b))$.
\end{proof}

So computing a greatest common divisor costs, up to a constant, the same as one schoolbook multiplication of $a$ by $b$. The
product of the quotients does the accounting: a step with a large quotient is expensive, but it shrinks the numbers a lot.
This is our first \emph{amortised} argument, and Section~\ref{sec:amortised} makes the idea systematic.

\section{Divide and conquer}\label{sec:master}

Many good algorithms cut a problem of size $n$ into $a$ problems of size $n/b$, solve them recursively, and combine the answers in time
$f(n)$. \emph{Merge sort} cuts a list in two halves, sorts each, and merges the two sorted halves by repeatedly moving the smaller of the
two front elements to the output: $a=b=2$, and merging two lists of total length $n$ takes at most $n-1$ comparisons. \emph{Binary
search} has $a=1$, $b=2$, $f(n)=1$. Karatsuba's multiplication (Chapter~9) has $a=3$, $b=2$, $f(n)=O(n)$; its story belongs here,
because it is the story of a lower bound that failed. In the autumn of 1960 Andrei Kolmogorov opened a seminar in Moscow on the complexity of
computation with the conjecture that multiplying two $n$-digit numbers needs about $n^2$ operations. Anatoly Karatsuba, a student of
twenty-three in the audience, found a method with about $n^{1.585}$ within a week; Kolmogorov announced it at the next meeting and closed
the seminar.

The running time satisfies $T(n)=aT(n/b)+f(n)$. Unrolling the recursion draws a tree (Figure~\ref{fig:rectree}): the root costs $f(n)$, its
$a$ children cost $f(n/b)$ each, level $j$ has $a^j$ nodes of cost $f(n/b^j)$, and after $\log_bn$ levels the $a^{\log_bn}=n^{\log_ba}$ leaves
cost $T(1)$ each. Everything depends on the competition between the root and the leaves.

\begin{figure}[ht]
\centering
\fig{ch01-rectree}
\caption{The recursion tree of $T(n)=3T(n/2)+f(n)$. Level $j$ costs $3^jf(n/2^j)$ in total. If the level totals decrease geometrically the
root wins and $T(n)=\Theta(f(n))$; if they increase geometrically the leaves win and $T(n)=\Theta(n^{\log_23})$; if they are all about equal,
$T(n)$ is the cost of a level times the number of levels.}
\label{fig:rectree}
\end{figure}

\begin{lemma}\label{lem:unroll}
Let $a>0$, $b\ge2$ an integer, and $T(b^K)=aT(b^{K-1})+f(b^K)$ for $K\ge1$. Then, for $n=b^K$,
\[
T(n)=n^{\log_ba}\,T(1)+\sum_{j=0}^{K-1}a^jf(n/b^j).
\]
\end{lemma}
\begin{proof}
By induction on $K$; for $K=0$ the sum is empty. For $K\ge1$, by induction applied to $n/b$, and since $a\cdot(n/b)^{\log_ba}=n^{\log_ba}$
(both sides have logarithm in base $b$ equal to $(K-1)\log_ba+\log_ba$),
\[
\begin{aligned}
T(n)&=a\Bigl((n/b)^{\log_ba}T(1)+\sum_{j=0}^{K-2}a^jf(n/b^{j+1})\Bigr)+f(n)\\
&=n^{\log_ba}T(1)+\sum_{j=1}^{K-1}a^jf(n/b^j)+f(n).
\end{aligned}
\]
\end{proof}

\begin{theorem}[Master theorem; Bentley, Haken and Saxe, 1980]\label{thm:master}
Let $a\ge1$, $b\ge2$ an integer, $f$ a positive function on the positive integers, and $T(1)>0$ and $T(n)=aT(n/b)+f(n)$ for $n$ a power
of $b$. Write $e=\log_ba$. Then, for $n$ running through the powers of $b$:
\begin{enumerate}
\item[(a)] if $f(n)=O(n^{e-\varepsilon})$ for some $\varepsilon>0$, then $T(n)=\Theta(n^e)$;
\item[(b)] if $f(n)=\Theta(n^e\log^kn)$ for some $k\ge0$, then $T(n)=\Theta(n^e\log^{k+1}n)$;
\item[(c)] if $f(n)=\Omega(n^{e+\varepsilon})$ for some $\varepsilon>0$, and there is $c<1$ with $af(n/b)\le cf(n)$ for all $n\ge b$, then
$T(n)=\Theta(f(n))$.
\end{enumerate}
\end{theorem}
\begin{proof}
Write $n=b^K$ and $S=\sum_{j=0}^{K-1}a^jf(n/b^j)$, so that $T(n)=n^eT(1)+S$ by Lemma~\ref{lem:unroll}. A bound that holds for all large
integers holds for all integers $m\ge1$ (or all $m\ge b$) after enlarging the constant, because only finitely many $m$ are left and
the functions involved are positive there; we use this silently. Note that $a^j(n/b^j)^e=a^jn^e/a^j=n^e$, since $(b^j)^e=a^j$.

(a) Now $f(m)\le Cm^{e-\varepsilon}$ for all $m\ge1$, so
\[
S\le C\sum_{j=0}^{K-1}a^j\Bigl(\frac n{b^j}\Bigr)^{e-\varepsilon}=Cn^{e-\varepsilon}\sum_{j=0}^{K-1}b^{\varepsilon j}=Cn^{e-\varepsilon}\frac{b^{\varepsilon K}-1}{b^\varepsilon-1}\le\frac{C}{b^\varepsilon-1}\,n^e,
\]
since $b^{\varepsilon K}=n^\varepsilon$. So $n^eT(1)\le T(n)\le(T(1)+C/(b^\varepsilon-1))n^e$.

(b) All the arguments $n/b^j$, $j<K$, are at least $b$, so $\log(n/b^j)>0$, and there are $0<c_1\le c_2$ with
$c_1m^e\log_b^km\le f(m)\le c_2m^e\log_b^km$ for all $m\ge b$ (the base of the logarithm only changes the constants). Then
\[
S=\Theta\Bigl(\sum_{j=0}^{K-1}n^e(K-j)^k\Bigr)=\Theta\Bigl(n^e\sum_{i=1}^Ki^k\Bigr).
\]
The sum $\sum_{i=1}^Ki^k$ is at most $K\cdot K^k$, and at least the sum of its terms with $i\ge K/2$, which is at least $(K/2)(K/2)^k$. So
$S=\Theta(n^eK^{k+1})=\Theta(n^e\log^{k+1}n)$, which dominates $n^eT(1)$ (once $K\ge1$).

(c) Iterating the regularity condition, $a^jf(n/b^j)\le c^jf(n)$ for $j\le K-1$: at each of the $j$ steps we apply it to some
$n/b^{i}$ with $i<j\le K-1$, which is at least $b$. So $f(n)\le S\le f(n)\sum_jc^j\le f(n)/(1-c)$. Finally
$n^eT(1)=O(n^{e})=O(f(n))$, since $f(n)\ge c'n^{e+\varepsilon}\ge c'n^e$. So $T(n)=\Theta(f(n))$.
\end{proof}

\begin{example}\label{exa:master}
Merge sort, $T(n)=2T(n/2)+\Theta(n)$: $e=1$, case (b) with $k=0$, $T(n)=\Theta(n\log n)$. Binary search, $T(n)=T(n/2)+1$: $e=0$, case (b),
$T(n)=\Theta(\log n)$. Karatsuba, $T(n)=3T(n/2)+\Theta(n)$: $e=\log_23\approx1.585$, case (a), $T(n)=\Theta(n^{\log_23})$. Strassen's matrix
product (Chapter~9), $T(n)=7T(n/2)+\Theta(n^2)$: case (a), $\Theta(n^{\log_27})=O(n^{2.81})$. And $T(n)=2T(n/2)+n^2$ is case (c), since
$2(n/2)^2=\frac12n^2$: $T(n)=\Theta(n^2)$, the root does most of the work.
\end{example}

The three cases do not cover everything: $f(n)=n^e/\log n$ falls between (a) and (b), and Exercise~\ref{ex:gap} shows that
$T(n)=\Theta(n^e\log\log n)$ there. More importantly, real algorithms cut $n$ into $\lceil n/2\rceil$ and $\lfloor n/2\rfloor$, not exactly in half.
The next result says that this rarely matters.

\begin{proposition}\label{prop:alln}
Let $T$ and $g$ be nondecreasing positive functions on the positive integers, and suppose $g(bm)\le Dg(m)$ for all $m$ and some constant
$D$. If $T(n)=\Theta(g(n))$ for $n$ running through the powers of $b$, then $T(n)=\Theta(g(n))$ for all $n$.
\end{proposition}
\begin{proof}
Let $c\,g(b^K)\le T(b^K)\le C\,g(b^K)$ for all $K$. Given $n$, take $K$ with $b^K\le n<b^{K+1}$. Since $T$ and $g$ are nondecreasing,
\[
T(n)\le T(b^{K+1})\le Cg(b^{K+1})\le CDg(b^K)\le CDg(n),
\]
\[
T(n)\ge T(b^K)\ge cg(b^K)\ge\tfrac cDg(b^{K+1})\ge\tfrac cDg(n).
\]
\end{proof}

The answers of cases (a) and (b), written $n^e(1+\log_bn)^k$ to keep them positive (this changes nothing up to constants), satisfy the
hypothesis with $D=b^e2^k$, since $1+\log_b(bm)\le2(1+\log_bm)$; so does $f$ in case (c) when $f$ is a power of $n$. For merge sort one can do better than
$\Theta$, and see exactly what ``about $n\log_2n$'' means.

\begin{proposition}\label{prop:mergesort}
Merge sort sorts $n$ elements with at most $n\lceil\log_2n\rceil$ comparisons.
\end{proposition}
\begin{proof}
Let $C(n)$ be the worst-case number of comparisons; $C(1)=0$ and $C(n)\le C(\lceil n/2\rceil)+C(\lfloor n/2\rfloor)+n-1$, since a merge removes
one element per comparison and the last element needs none. For $n\ge2$ we have $\lceil\log_2\lceil n/2\rceil\rceil=\lceil\log_2n\rceil-1$: if $2^{t-1}<n\le2^t$ then
$2^{t-2}<n/2\le\lceil n/2\rceil\le2^{t-1}$. So by induction, writing $t=\lceil\log_2n\rceil$,
\[
C(n)\le\lceil n/2\rceil(t-1)+\lfloor n/2\rfloor(t-1)+n-1=nt-1<nt.
\]
\end{proof}

When the pieces have different sizes the recursion tree is lopsided, but its level totals still tell the story.

\begin{proposition}\label{prop:twopieces}
Let $0<\alpha,\beta<1$, $c,C_0>0$, $n_0\ge2/\min(\alpha,\beta)$, and let $T\ge0$ on the positive integers satisfy $T(n)\le C_0$ for $n<n_0$ and
\[
T(n)\le T(\lfloor\alpha n\rfloor)+T(\lfloor\beta n\rfloor)+cn\qquad(n\ge n_0).
\]
\begin{enumerate}
\item[(i)] If $\alpha+\beta<1$, then $T(n)\le Dn$ for all $n$, with $D=\max\bigl(C_0,c/(1-\alpha-\beta)\bigr)$.
\item[(ii)] If $\alpha+\beta=1$, then $T(n)\le Dn\log_2n$ for all $n\ge2$, with $D=\max\bigl(C_0,c/h\bigr)$ and $h=-\alpha\log_2\alpha-\beta\log_2\beta>0$.
\end{enumerate}
\end{proposition}
\begin{proof}
For $n\ge n_0$ both $\lfloor\alpha n\rfloor$ and $\lfloor\beta n\rfloor$ are at least $2$ and less than $n$, so we may argue by induction on $n$.
(i) For $n<n_0$, $T(n)\le C_0\le Dn$. For $n\ge n_0$, $T(n)\le D\alpha n+D\beta n+cn=Dn-\bigl(D(1-\alpha-\beta)-c\bigr)n\le Dn$.
(ii) For $2\le n<n_0$, $T(n)\le C_0\le Dn\log_2n$. For $n\ge n_0$, since $x\mapsto x\log_2x$ is increasing for $x\ge1$ and $\lfloor\alpha n\rfloor\le\alpha n$,
\[
\begin{aligned}
T(n)&\le D\alpha n\log_2(\alpha n)+D\beta n\log_2(\beta n)+cn\\
&=Dn\log_2n-(Dh-c)n\le Dn\log_2n.
\end{aligned}
\]
Finally $h>0$ because $\log_2\alpha$ and $\log_2\beta$ are negative.
\end{proof}

So $T(n)=T(n/3)+T(2n/3)+n$ gives $O(n\log n)$ (here $h=\log_23-\frac23\approx0.918$), and the selection algorithm of Chapter~2, whose recurrence has the
shape $T(n/5)+T(7n/10)+cn$ with $\frac15+\frac7{10}<1$, runs in linear time. The quantity $h$ is the entropy of the split, which Tome~C4 will meet again.

\begin{remark}
The general statement is the \emph{Akra--Bazzi theorem} (1998): if
$T(n)=\sum_ia_iT(b_in)+f(n)$ with $0<b_i<1$, let $p$ solve $\sum_ia_ib_i^p=1$; then, under mild conditions on $f$,
$T(n)=\Theta\bigl(n^p(1+\int_1^nf(u)u^{-p-1}\,du)\bigr)$. Proposition~\ref{prop:twopieces} is the case of two pieces with $p=1$ (part (ii)) or
$p<1$ (part (i), where the integral is $O(n^{1-p})$).
\end{remark}

\section{Amortised analysis}\label{sec:amortised}

Some operations are usually cheap and occasionally expensive. A bank account is the right picture: cheap operations pay a little
more than they cost, and the savings pay for the expensive ones. Robert Tarjan made this into a method in 1985.

\begin{theorem}[Potential method]\label{thm:potential}
Let a sequence of operations take a data structure through states $D_0,D_1,\dots,D_m$, the $i$-th operation costing $c_i$. Let $\Phi$
be a real function of the states with $\Phi(D_i)\ge\Phi(D_0)$ for all $i$, and define the \emph{amortised cost} $\hat c_i=c_i+\Phi(D_i)-\Phi(D_{i-1})$.
Then $\sum_{i=1}^mc_i\le\sum_{i=1}^m\hat c_i$.
\end{theorem}
\begin{proof}
The sum telescopes: $\sum_i\hat c_i=\sum_ic_i+\Phi(D_m)-\Phi(D_0)\ge\sum_ic_i$.
\end{proof}

The proof is one line; the art is in choosing $\Phi$, the money in the bank. It should increase a little on cheap operations and fall a
lot on expensive ones.

\begin{proposition}[Binary counter]\label{prop:counter}
A $k$-bit binary counter starts at $0$. An increment flips bits from the right: the trailing $1$s become $0$, then the first $0$ (if any)
becomes $1$; its cost is the number of bits flipped. Then $m$ increments cost at most $2m$ in total.
\end{proposition}
\begin{proof}
Let $\Phi$ be the number of $1$s, so $\Phi\ge0=\Phi(D_0)$. An increment with $t$ trailing $1$s, followed by a $0$, costs $t+1$ and
changes $\Phi$ by $-t+1$: amortised cost $2$. If the counter is all $1$s, it costs $k$ and changes $\Phi$ by $-k$: amortised cost $0$.
Theorem~\ref{thm:potential} concludes.
\end{proof}

A single increment can cost $k$, and a naive bound would be $km$; the amortised bound says the average is below $2$. Figure~\ref{fig:potential} shows the
bank account at work: the actual total cost jumps at every power of $2$, when many bits flip at once, but it never crosses the line $2m$,
because the potential (the savings) is spent exactly then.

\begin{figure}[ht]
\centering
\fig{ch01-potential}
\caption{The first $40$ increments of a binary counter. Blue: the total cost so far. Green: the potential, the number of $1$s. Since
each increment has amortised cost exactly $2$, total cost plus potential equals $2m$, and the blue curve stays below the dashed line $2m$.}
\label{fig:potential}
\end{figure}
 Counting directly
gives the same: bit $i$ flips every $2^i$ increments, so the total is at most $\sum_{i\ge0}m/2^i=2m$. The potential is a way of doing such
counts when there is no clean formula.

\begin{proposition}[Dynamic array]\label{prop:dynarray}
An array grows by appending at the end. It starts with capacity $0$; when an append finds the array full (size $s$ equal to
capacity $c$) it allocates an array of capacity $\max(1,2c)$ and copies the $s$ elements over, at cost $s$; writing the new element
costs $1$. Then $m$ appends cost at most $3m$.
\end{proposition}
\begin{proof}
Take $\Phi=2s-c$. Initially $\Phi=0$. After the first append the capacity is a power of $2$ and $c\le 2s$ (a doubling happens at
$s=c$, giving capacity $2c$ for $c+1$ elements), so $\Phi\ge0$ throughout. An append without copying costs $1$ and raises $\Phi$ by $2$:
amortised $3$. An append that doubles a capacity $c\ge1$ costs $c+1$ and changes $\Phi$ from $2c-c=c$ to $2(c+1)-2c=2$: amortised
$c+1+2-c=3$. The very first append costs $1$ and brings $\Phi$ from $0$ to $1$: amortised $2$.
\end{proof}

This is how Python lists and C++ vectors grow (with factors a little different from $2$), and why appending is ``$O(1)$''. Doubling
is essential: growing by a fixed amount $d$ would cost about $m^2/(2d)$ (Exercise~\ref{ex:shrink} treats shrinking).

\begin{think}
The amortised bound is a statement about whole sequences. If you need every single operation to be fast (in a pacemaker, or
in a game that must draw a frame every $16$ milliseconds), a cost of $c$ once in a while may be unacceptable. Can you modify the dynamic
array so that every append costs $O(1)$, by copying a few old elements at each append instead of all at once?
\end{think}

\subsection*{Horn clauses in linear time}

Tome~C1 (Theorem~1.16) found the least model of a set of Horn clauses by repeatedly firing a rule $a_1\land\dots\land a_k\to b$ whose
hypotheses are all true. Searching for such a rule from scratch each time costs up to the length of the formula per firing, hence
quadratic time. Dowling and Gallier (1984) keep a counter.

\begin{theorem}[Dowling and Gallier, 1984]\label{thm:hornlinear}
The satisfiability of a set $S$ of Horn clauses can be decided, and its least model found, in time $O(L)$, where $L$ is the total
number of literal occurrences in $S$.
\end{theorem}
\begin{proof}
Prepare, in time $O(L)$: for each definite clause $C$, a counter $\mathrm{cnt}[C]$ equal to its number of hypothesis occurrences
(repetitions counted); for each variable $v$, the list $\mathrm{occ}[v]$ of the definite clauses in which $v$ occurs as a hypothesis, once
per occurrence; a queue $Q$ containing the conclusions of the facts (the clauses with $\mathrm{cnt}=0$); all variables false. Then
repeat while $Q$ is not empty: remove $v$ from $Q$; if $v$ is already true, do nothing; otherwise make $v$ true, and for each $C$ in
$\mathrm{occ}[v]$ decrease $\mathrm{cnt}[C]$, and if it reaches $0$ put the conclusion of $C$ into $Q$. Finally check each goal clause.

\emph{Every variable made true is true in every model of $S$}, by induction in the order of the queue: a conclusion enters $Q$ only when
its clause has all hypotheses true, and these are true in every model by induction.

\emph{The final assignment $M$ satisfies every definite clause.} If all the hypotheses of $C$ are true at the end, all its occurrences
were processed (each variable that becomes true has its whole list scanned), so $\mathrm{cnt}[C]$ reached $0$, its conclusion entered $Q$, and it was made true when removed.

So $M$ is the least model of the definite clauses of $S$, and as in Tome~C1, Theorem~1.16, $S$ is satisfiable if and only if every goal
clause has a false variable in $M$.

\emph{Time.} Each variable is made true at most once, so each list $\mathrm{occ}[v]$ is scanned at most once: $O(L)$ decrements in all.
A conclusion enters $Q$ at most once per clause (when its counter hits $0$, or initially for a fact), so the queue sees at most as many
insertions as there are clauses. The final check reads each goal once. Total $O(L)$.
\end{proof}

The aggregate argument (``each occurrence is decremented once'') is a potential argument in disguise: the potential is the sum of the
counters plus the number of entries in the queue.

\section{Lower bounds}\label{sec:lower}

An upper bound needs one good algorithm. A lower bound is a statement about \emph{every} algorithm, including all those nobody has
thought of yet. For general computation we are very bad at this (Tome~C3 explains why), but if the algorithm may only look at its input
through comparisons, there is a clean theory.

A \emph{comparison algorithm} on elements $x_1,\dots,x_n$ of a totally ordered set (assumed distinct) chooses, at each step, two indices $i,j$
and asks whether $x_i<x_j$; which question it asks next may depend on all previous answers, and the final output depends only on the
answers. Unrolled, the algorithm is a \emph{decision tree}: a binary tree whose inner nodes are questions $x_i:x_j$, whose edges are answers,
and whose leaves are outputs (Figure~\ref{fig:dtree}). The number of comparisons on an input is the depth of the leaf it reaches.

\begin{figure}[ht]
\centering
\fig{ch01-dtree}
\caption{The decision tree of insertion sort on three elements. The $3!=6$ orderings reach $6$ different leaves, so some leaf has depth
at least $\lceil\log_26\rceil=3$.}
\label{fig:dtree}
\end{figure}

\begin{lemma}\label{lem:leafdepth}
A binary tree with $L$ leaves has a leaf of depth at least $\log_2L$, and the average depth of its leaves is at least $\log_2L$.
\end{lemma}
\begin{proof}
A tree of depth $d$ has at most $2^d$ leaves, which gives the first claim. For the second, let $D(t)$ be the sum of the depths of the leaves of a
tree $t$ with $L$ leaves; we show $D(t)\ge L\log_2L$ by induction on the number of nodes. A single leaf has $D=0$. If the root has one
child, removing the root decreases every depth by $1$. If it has two subtrees with $L_1,L_2\ge1$ leaves, $L_1+L_2=L$, then each leaf is one
deeper than in its subtree, and by induction and the convexity of $x\mapsto x\log_2x$,
\[
D(t)\ge L+L_1\log_2L_1+L_2\log_2L_2\ge L+2\cdot\tfrac L2\log_2\tfrac L2=L\log_2L.
\]
\end{proof}

\begin{theorem}\label{thm:sortlower}
Every comparison algorithm that sorts $n$ distinct elements makes at least $\lceil\log_2n!\rceil$ comparisons on some input, and at least
$\log_2n!$ comparisons on average over the $n!$ orderings. Moreover $\log_2n!\ge n\log_2n-n\log_2e$.
\end{theorem}
\begin{proof}
Two different orderings of the input cannot reach the same leaf: the output, the same rearrangement in both cases, would be wrong for
one of them. So the tree has at least $n!$ leaves reached, and the lemma, applied to the tree pruned to the reached leaves, gives both bounds
(depths are integers, hence the ceiling). For the estimate, $e^n=\sum_k n^k/k!\ge n^n/n!$, so $n!\ge(n/e)^n$.
\end{proof}

With Proposition~\ref{prop:mergesort}: the best worst case is between $n\log_2n-1.443n$ and $n\log_2n$. Merge sort is optimal up to a
lower-order term. The exact minimum $S(n)$ is a famous puzzle. Lester Ford and Selmer Johnson found in 1959 a cunning method, \emph{merge
insertion}, which needs $\sum_{k=1}^n\lceil\log_2(3k/4)\rceil$ comparisons; this equals $\lceil\log_2 n!\rceil$ for $n\le11$, and for $n=12$ it is $30$
while $\lceil\log_212!\rceil=29$. In 1965 Mark Wells proved by a computer search that $S(12)=30$: the information bound is not always the
truth. Sorting five elements in $7=\lceil\log_2120\rceil$ comparisons is Exercise~\ref{ex:sort5}, and the notebook runs merge insertion.

\subsection*{Adversaries}

The information bound is weak for problems with few possible answers: finding the largest of $n$ elements has only $n$ answers, so
it gives $\log_2n$, while everybody knows that a knockout tournament needs $n-1$ matches. To prove such bounds we play against an
\emph{adversary}, who does not fix the input in advance but answers each question so as to keep the algorithm in the dark, while staying
consistent with \emph{some} input. The basic move is this.

\begin{lemma}\label{lem:raise}
Run a comparison algorithm on an input $x$. If an element $x_i$ never lost a comparison (it was never the smaller one), then raising $x_i$ above
all other elements gives an input on which the algorithm asks the same questions, receives the same answers and produces the same output.
The same holds for lowering an element that never won.
\end{lemma}
\begin{proof}
By induction on the steps: if the questions and answers agree so far, the next question is the same, since it depends only on previous
answers. Its answer is the same: a comparison not involving $x_i$ is unchanged, and one involving $x_i$ was won by $x_i$, which remains true
after raising $x_i$. The output depends only on the answers.
\end{proof}

\begin{proposition}\label{prop:max}
Every comparison algorithm that finds the largest of $n$ elements makes at least $n-1$ comparisons on every input.
\end{proposition}
\begin{proof}
If two elements never lost, raising either of them (Lemma~\ref{lem:raise}) gives two inputs with the same output and different maxima,
so the algorithm is wrong on one of them. Hence at most one element never lost: at least $n-1$ elements lost at least once. Each comparison
has one loser.
\end{proof}

\begin{theorem}[Pohl, 1972]\label{thm:maxmin}
Finding both the largest and the smallest of $n\ge2$ elements needs $\lceil3n/2\rceil-2$ comparisons in the worst case, and this many suffice.
\end{theorem}
\begin{proof}
\emph{Upper bound.} Compare the elements in pairs ($\lfloor n/2\rfloor$ comparisons), find the largest of the winners and the smallest of the losers,
putting a leftover element (for $n$ odd) in both groups. For $n$ even this is $n/2+2(n/2-1)=3n/2-2$; for $n$ odd, $(n-1)/2+2\cdot\frac{n-1}{2}=\frac{3n-3}{2}=\lceil3n/2\rceil-2$.

\emph{Lower bound.} Call an element \emph{fresh} if it has never been compared, and let $\Phi$ be the number of elements that never lost plus the
number that never won; a fresh element counts twice, so initially $\Phi=2n$. At the end, as in Proposition~\ref{prop:max}, at most one
element never lost and at most one never won, so $\Phi\le2$. The adversary keeps actual values, which it changes only by Lemma~\ref{lem:raise},
so its answers are always consistent. When asked to compare $x$ and $y$:
\begin{itemize}
\item if one of them never lost and the other has lost, it raises the one that never lost above everything and declares it the winner;
\item otherwise, if one never won and the other has won, it lowers the one that never won and declares it the loser;
\item otherwise it answers according to the current values.
\end{itemize}
In the first case the loser has lost before and keeps its status, and only the winner may stop being never-won; in the second case
the winner has won before and keeps its status, and only the loser may stop being never-lost. So $\Phi$ drops by at most $1$. In the third case, either both never lost and both never won (two fresh elements: $\Phi$ drops by $2$),
or both have lost (then the winner may stop being never-won; $\Phi$ drops by at most $1$), or both have won (symmetrically, at most $1$).
A comparison of two fresh elements removes two fresh elements, and no element becomes fresh again, so there are at most $\lfloor n/2\rfloor$ of them.
So the number $N$ of comparisons satisfies $2\lfloor n/2\rfloor+(N-\lfloor n/2\rfloor)\ge2n-2$, that is $N\ge2n-2-\lfloor n/2\rfloor=\lceil3n/2\rceil-2$.
\end{proof}

\subsection*{Lewis Carroll's tennis tournament}

In 1883 the Reverend Charles Dodgson, better known as Lewis Carroll, published a pamphlet on lawn-tennis tournaments, complaining that
the second prize went to the loser of the final, although the second-best player may have met the champion in the first round and gone
home. How many matches, then, does it take to find the best \emph{and} the second best? The champion has to win
about $\log_2n$ matches, and the second best is among the players the champion beat.

\begin{theorem}[Kislitsyn, 1964]\label{thm:second}
Finding the largest and the second largest of $n\ge2$ elements needs $n+\lceil\log_2n\rceil-2$ comparisons in the worst case, and this many
suffice.
\end{theorem}
\begin{proof}
\emph{Upper bound.} Play a knockout tournament in $\lceil\log_2n\rceil$ rounds, pairing the remaining players in each round (one gets a bye if
their number is odd): $n-1$ matches, since each eliminates one player. The champion played at most $\lceil\log_2n\rceil$ matches. The second
largest lost only to the champion, so it is the best of the at most $\lceil\log_2n\rceil$ players the champion beat: at most
$\lceil\log_2n\rceil-1$ more matches.

\emph{Lower bound.} Give every element a weight, initially $1$. The adversary keeps actual values, changed only by Lemma~\ref{lem:raise}. To
compare $x$ and $y$: if both never lost, the heavier wins (either one if equal), is raised above everything, and takes the weight of the
loser, whose weight becomes $0$; if exactly one never lost, it wins and is raised, weights unchanged; if both have lost, it answers by the
current values. The elements that never lost always carry the total weight $n$.

At the end at most one element never lost (Proposition~\ref{prop:max}), so it is the largest, $m$, and has weight $n$. Each time $m$ beat
an element that had never lost, its weight at most doubled, since the heavier one wins; its other wins did not change its weight. So $m$
beat at least $\lceil\log_2n\rceil$ elements $y_1,\dots,y_k$ that lost for the first time to $m$. Let $s$ be the output second largest, and suppose
some $y_i\ne s$ lost only once. Set the value of $y_i$ just below that of $m$ and both above all others: $m$ never lost and $y_i$ lost only
to $m$, so, as in Lemma~\ref{lem:raise}, all answers are unchanged, the output is still $s$, but now $y_i$ is the second largest. So all
the $y_i$ but at most one lost at least twice. Counting losers, one per comparison: $n-1$ elements lost at least once, and at least $k-1$
of them lost a second time, so there are at least $n-1+\lceil\log_2n\rceil-1$ comparisons.
\end{proof}

So Carroll was right to worry, and the fix is cheap: after the final, about $\log_2n$ extra matches among the champion's victims find the true
runner-up. For $n=64$ players that is $68$ matches in all instead of the $63$ of the tournament, and nobody can do better.

\begin{beauty}
A lower bound is a theorem about algorithms that do not exist yet. For comparisons there are two weapons: \emph{information}, which counts
the answers ($n!$ orders need $\log_2n!$ yes-or-no questions), and the \emph{adversary}, who answers so as to reveal as little as possible.
Lam\'e's question about Euclid, Carroll's complaint about tennis and Kolmogorov's failed conjecture about multiplication are the same
question asked three times: how fast can it possibly be done? Sometimes, as for sorting and for the second best, the answer comes with a
matching algorithm; sometimes, as for multiplication, the conjecture falls and the search goes on.
\end{beauty}

\section{Exercises}

\begin{exercise}[$\star$]\label{ex:recsolve}
Solve, for $n$ a power of the base: $T(n)=4T(n/2)+n^2$; $T(n)=9T(n/3)+n$; $T(n)=2T(n/4)+\sqrt n$; $T(n)=T(n/2)+n$.
\end{exercise}

\begin{exercise}[$\star$]\label{ex:gap}
Let $T(1)=1$ and $T(n)=2T(n/2)+n/\log_2n$ for $n=2^K$, $K\ge1$. Show that $T(n)=\Theta(n\log\log n)$, and conclude that the master theorem
cannot be extended to cover $f(n)=n^e/\log n$ with answer $\Theta(n^e)$ or $\Theta(n^e\log n)$.
\end{exercise}

\begin{exercise}[$\star\star$]\label{ex:stooge}
\emph{Stooge sort} sorts $A[1..n]$ as follows: if $n\le2$, sort directly; otherwise, with $k=\lceil2n/3\rceil$, sort the first $k$ elements
recursively, then the last $k$, then the first $k$ again. Prove that it sorts, and find its running time.
\end{exercise}

\begin{exercise}[$\star$]\label{ex:multipop}
A stack supports $\mathrm{push}(x)$ (cost $1$), $\mathrm{pop}$ (cost $1$) and $\mathrm{multipop}(k)$, which pops $\min(k,s)$ elements where $s$ is the
current size (cost $\min(k,s)$). Show that any sequence of $m$ operations on an initially empty stack costs at most $2m$.
\end{exercise}

\begin{exercise}[$\star$]\label{ex:twostacks}
Implement a queue with two stacks: enqueue pushes on the stack $\mathrm{In}$; dequeue pops from $\mathrm{Out}$, first moving all of
$\mathrm{In}$ onto $\mathrm{Out}$ (one pop and one push per element) if $\mathrm{Out}$ is empty. Show that $m$ operations cost $O(m)$ stack
operations, although a single dequeue may cost $\Theta(m)$.
\end{exercise}

\begin{exercise}[$\star\star$]\label{ex:shrink}
A dynamic array also supports deleting the last element. (a) Show that the rule ``halve the capacity when the array is half full'' can make
$m$ operations cost $\Theta(m^2)$. (b) With the rule ``when the size drops to a quarter of the capacity $c\ge4$, copy into an array of capacity
$c/2$'', show that $m$ operations cost $O(m)$, using $\Phi=2s-c$ when $s\ge c/2$ and $\Phi=c/2-s$ when $s<c/2$.
\end{exercise}

\begin{exercise}[$\star\star$]\label{ex:merge}
Show that merging two sorted lists $a_1<\dots<a_m$ and $b_1<\dots<b_m$ by comparisons needs $2m-1$ comparisons in the worst case. Compare with
the information bound $\log_2\binom{2m}m$.
\end{exercise}

\begin{exercise}[$\star\star\star$]\label{ex:sort5}
Sort five elements with $7$ comparisons in the worst case, which is optimal.
\end{exercise}

\begin{exercise}[$\star$]\label{ex:insertion}
Insertion sort inserts $x_2,x_3,\dots,x_n$ in turn into the sorted list of the previous ones, scanning from the right. Show that the number of
comparisons is between $I$ and $I+n-1$, where $I$ is the number of \emph{inversions} (pairs $i<j$ with $x_i>x_j$), and that on a uniformly random
ordering $\mathbb E[I]=n(n-1)/4$.
\end{exercise}

\begin{exercise}[$\star$]\label{ex:fibcalls}
The naive recursive program $\mathrm{fib}(n)=\mathrm{fib}(n-1)+\mathrm{fib}(n-2)$, with $\mathrm{fib}(0)=0$, $\mathrm{fib}(1)=1$, makes $c(n)$ calls in all. Show
$c(n)=2F_{n+1}-1$, so the running time grows like $\varphi^n$. Then compute $F_n$ with $O(\log n)$ multiplications of $2\times2$ matrices, using
$\left(\begin{smallmatrix}1&1\\1&0\end{smallmatrix}\right)^n=\left(\begin{smallmatrix}F_{n+1}&F_n\\F_n&F_{n-1}\end{smallmatrix}\right)$.
\end{exercise}

\section{Hints}

\begin{hint}{ex:recsolve}
Compute $e=\log_ba$ each time and compare $f$ with $n^e$; for the last one check the regularity condition.
\end{hint}
\begin{hint}{ex:gap}
Use Lemma~\ref{lem:unroll}: the level totals are $n/(K-j)$, and $\sum_{i=1}^K1/i$ is a harmonic sum.
\end{hint}
\begin{hint}{ex:stooge}
After the first two recursive calls, the largest $n-k$ elements are in their final places. Why does the overlap of the two thirds matter?
\end{hint}
\begin{hint}{ex:multipop}
Each element is pushed once and popped at most once; or take $\Phi$ the size of the stack.
\end{hint}
\begin{hint}{ex:twostacks}
Take $\Phi=2|\mathrm{In}|$.
\end{hint}
\begin{hint}{ex:shrink}
(a) Alternate insertions and deletions at the threshold. (b) Check that a resize happens only when $\Phi$ equals the number of elements copied.
\end{hint}
\begin{hint}{ex:merge}
Consider the input $a_1<b_1<a_2<b_2<\dots<a_m<b_m$: which comparisons must be made?
\end{hint}
\begin{hint}{ex:sort5}
Compare $a:b$ and $c:d$, then the two winners. Insert the fifth element into the chain of three, and finally the remaining element into
the elements known to be below its partner.
\end{hint}
\begin{hint}{ex:insertion}
Each comparison that moves $x_j$ one place to the left removes one inversion. For the average, use linearity of expectation over the pairs.
\end{hint}
\begin{hint}{ex:fibcalls}
Show $c(n)+1$ satisfies the Fibonacci recurrence. For the matrix power, square repeatedly.
\end{hint}

\section{Solutions}

\begin{solution}{ex:recsolve}
\idea Each time, the question is who wins: the root $f(n)$, the leaves $n^{\log_ba}$, or neither.
\writeup $4T(n/2)+n^2$: $e=2$, $f=\Theta(n^e)$, case (b), $\Theta(n^2\log n)$. $9T(n/3)+n$: $e=2$, $f=O(n^{2-1})$, case (a), $\Theta(n^2)$.
$2T(n/4)+\sqrt n$: $e=\frac12$, case (b), $\Theta(\sqrt n\log n)$. $T(n/2)+n$: $e=0$, $f(n)=n=\Omega(n^{0+1})$ and $f(n/2)=\frac12f(n)$, so case (c) with
$c=\frac12$, $\Theta(n)$.
\end{solution}

\begin{solution}{ex:gap}
\idea Between cases (a) and (b) the level totals are neither geometric nor constant, they are harmonic.
\writeup By Lemma~\ref{lem:unroll} with $n=2^K$, $T(n)=n+\sum_{j=0}^{K-1}2^j\frac{n/2^j}{K-j}=n+n\sum_{i=1}^K\frac1i=n(1+H_K)$. Since $\ln K\le H_K\le1+\ln K$ (compare with
$\int dx/x$) and $K=\log_2n$, $T(n)=\Theta(n\log\log n)$. This is neither $\Theta(n)$ nor $\Theta(n\log n)$, so no formula of the two shapes holds
for $f=n^e/\log n$.
\end{solution}

\begin{solution}{ex:stooge}
\idea Track where the largest $n-k$ elements go: the overlap of the two thirds is big enough to catch all of them.
\writeup By induction on $n$; the cases $n\le2$ are clear. Let $n\ge3$, $k=\lceil2n/3\rceil$, and let $B$ be the set of the $n-k$ largest elements.
Note $n-k\le k-(n-k)$, that is $2(n-k)\le k$, since $k\ge2n/3$. After the first call, the first $k$ positions are sorted, so every element of $B$
that was among them is among the last $n-k$ of them, at positions $>k-(n-k)$ (if an element of $B$ sat at a position $\le 2k-n$, the $k-(2k-n)=n-k$
positions after it, all within the first $k$, would hold larger elements, all in $B$, too many). Hence all of $B$ is in positions $>2k-n\ge n-k$,
which lie in the last $k$ positions. The second call sorts the last $k$ positions, putting $B$, sorted, at positions $k+1,\dots,n$, its final place.
The third call sorts the first $k$ positions, which contain the rest. Time: $T(n)=3T(2n/3)+O(1)$, so $e=\log_{3/2}3=\ln3/\ln1.5\approx2.71$ and
$T(n)=\Theta(n^{2.71})$, worse than insertion sort. (Here $b=3/2$ is not an integer; we assumed $b$ integral only so that $n/b^j$ is an integer, and ignoring rounding the proofs of
Lemma~\ref{lem:unroll} and of case (a) go through unchanged.)
\end{solution}

\begin{solution}{ex:multipop}
\idea An element can only be popped if it was pushed; pay for its pop when it is pushed.
\writeup Let $\Phi$ be the size of the stack, $\Phi\ge0=\Phi(D_0)$. A push costs $1$ and raises $\Phi$ by $1$: amortised $2$. A pop or a multipop
popping $p$ elements costs $p$ and lowers $\Phi$ by $p$: amortised $0$. By Theorem~\ref{thm:potential}, the total is at most $2m$.
\end{solution}

\begin{solution}{ex:twostacks}
\idea Each element is moved from $\mathrm{In}$ to $\mathrm{Out}$ at most once; prepay the move at enqueue.
\writeup Take $\Phi=2|\mathrm{In}|\ge0$. Enqueue: one push, $\Phi$ up by $2$: amortised $3$. Dequeue with $\mathrm{Out}$ non-empty: one pop, amortised $1$.
Dequeue with $\mathrm{Out}$ empty and $|\mathrm{In}|=t$: $2t$ moves plus one pop, $\Phi$ down by $2t$: amortised $1$. So $m$ operations cost at most
$3m$. A single dequeue after $m-1$ enqueues costs $2(m-1)+1$.
\end{solution}

\begin{solution}{ex:shrink}
\idea A bad rule lets the size oscillate around a resize threshold; a good one leaves room on both sides after every resize.
\writeup (a) Fill the array to size $s=c=2^t$ and then alternate: append (the array is full: copy $2^t$ elements into capacity $2^{t+1}$), delete
(now half full: copy $2^t$ elements back). Each pair costs about $2^{t+1}$; with $2^t\approx m/2$ initial appends and $m/2$ further operations the
total is $\Theta(m^2)$. (b) $\Phi\ge0$ in both regimes and $\Phi=0$ initially. Check the four kinds of operations. Append with $s\ge c/2$, no
resize: cost $1$, $\Phi$ up by $2$, amortised $3$. Append with $s<c/2$: $\Phi$ goes down by $1$, amortised $0$. Append with $s=c$: doubling as in
Proposition~\ref{prop:dynarray}, amortised $3$. Delete with $s>c/2$: cost $1$, $\Phi$ down by $2$, amortised $\le0$. Delete with $s\le c/2$, no resize:
cost $1$, $\Phi$ up by $1$, amortised $2$. Delete triggering a shrink: before the deletion $s=c/4+1$, $\Phi=c/4-1$; it costs $1$ plus copying $c/4$ elements,
and afterwards $s=c/4=c'/2$ for the new capacity $c'=c/2$, so $\Phi=0$: amortised $1+c/4-(c/4-1)=2$. So every operation has amortised cost at most $3$.
\end{solution}

\begin{solution}{ex:merge}
\idea Find an input where every adjacent pair of the merged order is a pair from different lists, and argue each must be compared directly.
\writeup Upper bound: the usual merge uses at most $2m-1$. Lower bound: on the input $a_1<b_1<a_2<b_2<\dots<a_m<b_m$, consider two consecutive elements
$u<v$ of the merged order; they come from different lists. If the algorithm never compares $u$ with $v$, swap their values: $u$ and $v$ keep
their order relative to every other element (no element lies between them) and each list stays sorted (their neighbours in their own lists
lie outside the interval). So all answers are the same, the output is the same, and it is wrong for one of the two inputs. Hence all $2m-1$
consecutive pairs must be compared. The information bound $\log_2\binom{2m}{m}=2m-\frac12\log_2m+O(1)$ is weaker: the adversary does better than
counting.
\end{solution}

\begin{solution}{ex:sort5}
\idea Spend the first comparisons building a chain of three with a fourth element hanging below its top; then binary insertion into $3$ elements
costs only $2$ comparisons, and so does the last element.
\writeup Compare $a:b$ and $c:d$; rename so that $a<b$, $c<d$, and compare $b:d$, renaming the pairs so that $b<d$. Now $a<b<d$ and $c<d$ ($3$ comparisons).
Insert $e$ into the chain $a<b<d$ by binary search: compare with $b$, then with $a$ or $d$ ($2$ comparisons). We now have a chain of $4$ containing
$a,b,d,e$, and $c<d$, so $c$ must go among the elements of the chain below $d$: these are $a,b$ and possibly $e$, at most three, so binary
insertion takes $2$ comparisons. Total $7$, and $7=\lceil\log_2120\rceil$ is optimal by Theorem~\ref{thm:sortlower}.
\end{solution}

\begin{solution}{ex:insertion}
\idea Insertion sort repairs inversions one at a time, and each pair is inverted with probability one half.
\writeup When $x_j$ is inserted, it is compared with the previous elements larger than it, one by one, and then with one smaller element, unless
it reaches the front. The larger previous elements are exactly the inversions $(i,j)$ with $i<j$. So the comparisons for $x_j$ number between
$I_j$ and $I_j+1$, where $I_j$ is the number of inversions ending at $j$, and summing over $j=2,\dots,n$ gives $I\le\text{comparisons}\le I+n-1$. For a
random ordering each of the $\binom n2$ pairs is inverted with probability $\frac12$ by symmetry, so $\mathbb E[I]=\frac12\binom n2=\frac{n(n-1)}4$.
\end{solution}

\begin{solution}{ex:fibcalls}
\idea The number of calls satisfies the same recurrence as the Fibonacci numbers, up to $1$.
\writeup $c(0)=c(1)=1$ and $c(n)=1+c(n-1)+c(n-2)$. Then $d(n)=c(n)+1$ satisfies $d(n)=d(n-1)+d(n-2)$, $d(0)=d(1)=2=2F_1=2F_2$, so $d(n)=2F_{n+1}$ and
$c(n)=2F_{n+1}-1\approx2\varphi^{n+1}/\sqrt5$. The matrix identity holds for $n=1$ and passes from $n$ to $n+1$ by multiplying by
$\left(\begin{smallmatrix}1&1\\1&0\end{smallmatrix}\right)$. To compute $M^n$, write $M^n=(M^{\lfloor n/2\rfloor})^2$ or $(M^{\lfloor n/2\rfloor})^2M$: at most
$2\log_2n$ matrix products, each $8$ multiplications. (On numbers with $\Theta(n)$ bits, each multiplication is itself not one step; Chapter~9
counts the bits.)
\end{solution}

\section*{Sources and further reading}

\begin{itemize}
\item \emph{The stories.} G.~Lam\'e, \emph{Note sur la limite du nombre des divisions dans la recherche du plus grand commun diviseur entre deux nombres
entiers}, Comptes Rendus de l'Acad\'emie des Sciences 19 (1844), 867--870; J.~Shallit, \emph{Origins of the analysis of the Euclidean algorithm},
Historia Mathematica 21 (1994), 401--419, which describes the earlier bounds (among them P.-J.-\'E.~Finck's of 1841); for Lam\'e's life, the MacTutor
biography by J.~J.~O'Connor and E.~F.~Robertson (University of St Andrews).
The Knuth quotation ends his memo \emph{Notes on the van Emde Boas construction of priority deques: an instructive use of recursion} (1977), as
Knuth recalls in the frequently asked questions on his Stanford web page. A.~A.~Karatsuba, \emph{The complexity of computations},
Proceedings of the Steklov Institute of Mathematics 211 (1995), 169--183, tells the story of the 1960 seminar; the paper A.~Karatsuba and
Yu.~Ofman, \emph{Multiplication of multidigit numbers on automata}, Doklady Akademii Nauk SSSR 145 (1962), 293--294, was written by Kolmogorov.
C.~L.~Dodgson, \emph{Lawn Tennis Tournaments: The True Method of Assigning Prizes with a Proof of the Fallacy of the Present Method}, Macmillan, 1883;
the story and Kislitsyn's theorem are in D.~E.~Knuth, \emph{The Art of Computer Programming}, vol.~3, \emph{Sorting and Searching}, 2nd ed.,
Addison-Wesley, 1998, Section~5.3.3.
\item J.~Hartmanis and J.~Simon, \emph{On the power of multiplication in random access machines}, Proceedings of the 15th IEEE Symposium on
Switching and Automata Theory (1974), 13--23. For the word RAM, M.~L.~Fredman and D.~E.~Willard, \emph{Surpassing the information theoretic bound
with fusion trees}, Journal of Computer and System Sciences 47 (1993), 424--436.
\item J.~L.~Bentley, D.~Haken and J.~B.~Saxe, \emph{A general method for solving divide-and-conquer recurrences}, SIGACT News 12(3) (1980), 36--44;
M.~Akra and L.~Bazzi, \emph{On the solution of linear recurrence equations}, Computational Optimization and Applications 10 (1998), 195--210;
T.~H.~Cormen, C.~E.~Leiserson, R.~L.~Rivest and C.~Stein, \emph{Introduction to Algorithms}, 4th ed., MIT Press, 2022, Chapters~4 and~16, the
source of our presentation of the master theorem and of the dynamic array.
\item R.~E.~Tarjan, \emph{Amortized computational complexity}, SIAM Journal on Algebraic and Discrete Methods 6 (1985), 306--318;
W.~F.~Dowling and J.~H.~Gallier, \emph{Linear-time algorithms for testing the satisfiability of propositional Horn formulae}, Journal of Logic
Programming 1 (1984), 267--284.
\item L.~R.~Ford and S.~M.~Johnson, \emph{A tournament problem}, American Mathematical Monthly 66 (1959), 387--389; the value $S(12)=30$ (M.~Wells, 1965)
and the later history are in Knuth's volume~3, Section~5.3.1; G.~K.~Manacher, \emph{The Ford--Johnson sorting algorithm is not optimal},
Journal of the ACM 26 (1979), 441--456. I.~Pohl, \emph{A sorting problem and its complexity}, Communications of the ACM 15 (1972), 462--464;
S.~S.~Kislitsyn, \emph{On the selection of the $k$th element of an ordered set by pairwise comparisons}, Sibirskii Matematicheskii Zhurnal 5 (1964),
557--564.
\end{itemize}
